% Generated by paper-covert from output/manuscript/sections_revised/02_introduction.md

\section{Introduction}\label{introduction}

Australia's grains sector --- wheat, barley, canola, and pulse/legume crops --- is a major export industry, yet no continent-wide, field-verified, species-resolved crop map exists at paddock scale. Official production and area statistics from the Australian Bureau of Statistics (ABS) and the Australian Bureau of Agricultural and Resource Economics and Sciences (ABARES) are available only at coarse administrative-region granularity, with no spatial detail below the Statistical Area Level 2 (SA2) or state level, and the two agencies disagree with each other by a median 6.8\% on national sown area --- an honest floor against which any finer-grained map should be judged, not a target of zero disagreement. A spatially explicit, species-resolved product would let growers, agronomists, and policy analysts see where each crop is actually grown, at the scale a single farm operates, rather than inferring it from a regional aggregate.

Remote sensing has made continent- and global-scale cropland mapping technically feasible, but scale and species resolution trade off against each other in practice. Products built for global coverage typically classify land as cropland versus non-cropland, or into a small number of broad crop-type groups defined by growing calendar rather than species, because building species-resolved training data at global scale is prohibitively expensive. Australia sits in a Southern Hemisphere, canola-and-legume-growing agricultural system that most global crop-calendar-based products were not designed around, so even where global coverage exists, its species resolution and seasonal assumptions may not transfer.

Three threads of prior work bear on this problem. Global crop-mapping infrastructure --- WorldCereal \citep{vantricht2023worldcereal}, foundation-model embeddings such as Presto \citep{tseng2024presto} and AlphaEarth \citep{brown2025alphaearth}, and generic land-cover products such as Esri/Impact Observatory's global land-use-land-cover layer \citep{karra2021esri} --- delivers broad coverage but coarse or generic crop-type resolution, and none has been evaluated for Australian species-level generalizability. Two Australian precedents come closer: an in-season Sentinel-2 classifier for six classes in Western Australia \citep{sharma2026wa} and a Sentinel-1/Sentinel-2/MODIS fusion classifier distinguishing cereals from canola in the Murray-Darling Basin \citep{alshammari2024mdb}. Both are regional rather than continent-wide, and both train on either another model's output or opportunistic harvester-derived labels rather than field-verified ground truth. A third thread addresses the paddock-boundary problem this map also depends on: CSIRO's ePaddocks/Graincast system \citep{lawes2022graincast} is the closest operational Australian precedent for boundary-plus-yield mapping at scale, while the Segment Anything Model (SAM; \citealp{kirillov2023sam}) --- the segmentation backbone used here --- and Fields of the World \citep{kerner2024fields}, a cloud-robust Sentinel-1/Sentinel-2 boundary-delineation benchmark that does not cover Australia, represent the general-purpose segmentation methods this paper draws on instead of a bespoke Australian boundary product.

The gap this leaves is specific: no published map combines continent-wide coverage, paddock-scale resolution, species-level (not merely cropland) classification, and training exclusively on field-verified ground truth, for the Australian grains system. Closing it exactly as originally scoped --- a full \textasciitilde10-species classifier trained via foundation-model fine-tuning --- proved not to be reachable at the field-verified label volume available: a 9-species classifier reached only macro F1 0.38, and fine-tuned Presto embeddings scored worse than hand-built spectral indices (macro F1 0.775 vs.~0.821) at this label volume. We report both as genuine negative results rather than silently narrowing scope without comment, and instead ship a 3-group classifier (Canola, Cereal, Legume) built on hand-engineered, day-of-year-binned spectral indices.

This paper makes two primary contributions. First, we present a national, polygon-level, three-class crop-species map of Australia at 10 m resolution, built end-to-end from open Sentinel-2 imagery and 2,194 field-verified GRDC National Variety Trial (NVT) records --- not another model's predictions or opportunistic harvester data. Every polygon carries a predicted class or an explicit abstain reason, a per-polygon confidence, and, for Cereal, a calibrated yield estimate; the map is validated against independent ABS sown-area statistics with no connection to the training pipeline. Second, we quantify and mechanistically explain a systematic area-inflation failure mode common to presence-only crop mapping --- the map over-calls crop-present area by 1.47-1.59x, concentrated almost entirely in the Cereal and Legume classes and essentially absent from Canola --- and report a negative result from the natural, more principled fix: a phenology-shape presence gate that closed the area-inflation gap but did so by disproportionately rejecting real canola paddocks, and was rejected for that reason. This is a transferable caution for presence-gate design in crop mapping generally, not only for this pipeline. Two further findings support these primary contributions rather than standing alone: a label-conflict resolution specific to point-based agronomic trial networks, where collapsing species to broad groups and hand-reviewing co-located trials measurably improves classifier performance; and a negative result for foundation-model embeddings against hand-built spectral indices at this label volume, reported above.

The remainder of this paper is organized as follows. Section 2 reviews related work in global crop mapping, Australian precedents, paddock-boundary delineation, and accuracy-assessment standards. Section 3 describes the study area and data sources. Section 4 details the segmentation, labelling, classification, presence-gating, and yield-calibration methodology. Section 5 reports classifier accuracy, the ABS validation, the phenology-gate experiment, yield calibration, and an independent spatial comparison against the WorldCereal and National Land Use Map (NLUM) products. Section 6 discusses the area-inflation finding and its implications for presence-gate design. Section 7 states limitations, including the map's largest current scope gap --- single-season (2024) national coverage, pending a planned multi-year national extension --- and Section 8 concludes.
